% Part: methods
% Chapter: induction
% Section: strong-induction

\documentclass[../../../include/open-logic-section]{subfiles}

\begin{document}

\olfileid{mth}{ind}{str}

\olsection{استقرای قوی}

در اصل استقرایی که پیش‌تر بررسی شد، $P(0)$ را اثبات می‌کنیم و نیز
نشان می‌دهیم اگر $P(k)$، آنگاه $P(k+1)$. در بخش دوم، صدق $P(k)$ را
فرض و از این فرض برای اثبات $P(k+1)$ استفاده می‌کنیم. البته به‌طور
هم‌ارز می‌توانیم برای هر $k>0$، $P(k-1)$ را فرض کنیم و آن را برای
اثبات $P(k)$ به کار ببریم. بخش مهم این است که بتوانیم استنتاج از هر
عدد به جانشین آن را انجام دهیم؛ یعنی بتوانیم گزارهٔ مورد نظر را برای
هر عدد مثبت با فرض برقرار‌بودنش برای پیشین آن اثبات کنیم.

گونه‌ای دیگر از اصل استقرا وجود دارد که در آن فقط برقرار‌بودن گزاره
برای پیشین $k-1$ از $k$ را فرض نمی‌کنیم، بلکه برقرار‌بودن آن را برای
همهٔ اعداد کوچک‌تر از~$k$ فرض می‌کنیم و از این فرض برای اثبات گزاره
دربارهٔ~$k$ استفاده می‌کنیم. این روش نیز $P(n)$ را برای همهٔ
$n \in \Nat$ به ما می‌دهد. زیرا وقتی $P(0)$ را برقرار کرده باشیم،
در نتیجه نشان داده‌ایم $P$ برای همهٔ اعداد کوچک‌تر از~$1$ برقرار است.
اگر بدانیم «چنانچه $P(l)$ برای همهٔ $l<k$ برقرار باشد، آنگاه $P(k)$»،
این را به‌ویژه برای $k=1$ می‌دانیم؛ پس می‌توانیم $P(1)$ را نتیجه
بگیریم. بدین ترتیب $P(0)$ و $P(1)$، یعنی $P(l)$ برای همهٔ $l<2$، را
اثبات کرده‌ایم. چون این گزارهٔ شرطی را نیز داریم که اگر $P(l)$ برای
همهٔ $l<2$، آنگاه $P(2)$، می‌توانیم~$P(2)$ را نتیجه بگیریم، و
همین‌طور ادامه دهیم.

در واقع، اگر بتوانیم گزارهٔ شرطی کلی «برای هر $k$، اگر $P(l)$ برای
همهٔ $l<k$، آنگاه $P(k)$» را برقرار کنیم، دیگر لازم نیست $P(0)$ را
جداگانه اثبات کنیم، زیرا از همین گزاره نتیجه می‌شود. به یاد آورید که
گزاره‌ای کلی مانند «برای همهٔ $l<k$، $P(l)$» هنگامی که هیچ $l<k$ای
وجود ندارد صادق است. این نمونه‌ای از سور تهی است: «همهٔ $A$ها $B$
هستند» وقتی هیچ $A$ای وجود نداشته باشد صادق است، و
$\lforall[x][(!A(x) \lif !B(x))]$ وقتی هیچ $x$ای شرط~$!A(x)$ را برآورده
نکند صادق است. در اینجا صورت صوری‌شده چنین خواهد بود:
«$\lforall[l][(l < k \lif P(l))]$»؛ و اگر هیچ $l < k$ای وجود نداشته
باشد، این گزاره صادق است. برای $k=0$ دقیقاً چنین است: هیچ $l<0$ای وجود
ندارد؛ پس «برای همهٔ $l<0$، $P(l)$» هرچه $P$ باشد صادق است. بنابراین
اثبات «اگر $P(l)$ برای همهٔ $l<k$، آنگاه $P(k)$» به‌طور خودکار
~$P(0)$ را نیز برقرار می‌کند.

این گونه هنگامی سودمند است که اثبات گزاره دربارهٔ~$k$ را نتوان تنها
بر گزارهٔ مربوط به $k-1$ استوار کرد، بلکه شاید لازم باشد فرض کنیم آن
گزاره برای یک یا چند $l<k$ صادق است.

\end{document}
